\section{From prices to flags}\label{sec:flags}

Theorem~\ref{ext:star} asks for a matrix whose columns nearly attain
their combined star bound. At the tip, a bias threshold selects ordinary
subspaces, and the dimension of their span records how much of the
ambient space they supply. Away from the tip the same measurement,
obtained from the exterior price, can take fractional values. We need
to estimate it for many columns at once.

We first recover a dimension for sums, intersections, and images of
formal subspaces from these threshold measurements. The subspaces for
one site, ordered by bias, then form a flag. Intersecting the levels of
two flags gives them a pair law. Successively combining pairs of flags
will provide a count of the columns available at a threshold.
Lemma~\ref{flag:insertion} gives a lower bound on the rank of
Vandermonde columns by the expected count truncated at the number of
rows, up to arbitrarily small error. One matrix realizes this bound
at any prescribed finite set of thresholds.

\subsection{Threshold rank}

Let $Y$ be a Hadamard cone with tip $0$, and let $K$ be an infinite field.
A \emph{supplier} in a finite-dimensional $K$-space $E$ is a triple
$(A,y,\theta)$ with $A\subseteq E$ a subspace, $y\in Y$ a site and
$\theta\in\mathbb R$ a bias. Call $(y,\theta)$ its label. A list spans if
its supplied subspaces span $E$.
The exterior value in~\eqref{ext:value} assigns a finite number
$\lambda_E(L)$ to each spanning list $L$. In this section we may use
any such value with the following properties.
\begin{enumerate}
\item Isomorphisms of $E$ and permutations of the list preserve value.
The value is continuous in the labels, positively homogeneous under their
common dilation, and convex under simultaneous geodesic interpolation
of sites and linear interpolation of biases.
\item Increasing a bias, enlarging a subspace or adding a supplier can
only increase value. Shifting every bias by $t$ adds $t\dim E$.
Suppliers with the same site and bias merge by summing their subspaces,
and supplying the zero space has no effect.
\item Taking the direct sum of two spaces with their respective suppliers
adds their values. Adding a supplier has diminishing increments under
enlargement of a spanning base list.
\end{enumerate}
Section~\ref{sec:exterior} established these properties for the exterior
value. A supplier is redundant if another supplier at the same site
offers a containing subspace at a higher or equal bias. Indeed, raising
the first bias to the second allows the two suppliers to merge, and
monotonicity gives equality at the original bias.

For any list $L$ in $E$, append the \emph{root supplier} $(E,0,s)$ and
write $\lambda_L(s)$ for the value of the resulting spanning list. The
root offers all of $E$. Its marginal value measures how much dimension
it supplies at bias $s$. The remainder is the \emph{threshold rank}
of the ordinary suppliers. Writing $L-s$ for the list with $s$
subtracted from every ordinary bias, we have
\begin{equation}\label{flag:capdef}
 r_E(L)=\dim E-\lambda'_{L,+}(0),\qquad
 r_E(L-s)=\dim E-\lambda'_{L,+}(s).
\end{equation}
The second identity follows by translating all biases. Convexity
ensures the existence of the right derivative. Monotonicity and bias
translation bound it between zero and $\dim E$.

\begin{lemma}[Threshold rank]\label{flag:capacity}
The rank $r_E$ takes values in $[0,\dim E]$, adds on direct sums, and is
unchanged on merging equal labels or adding unused ambient dimensions.
It is monotone under adding suppliers, enlarging subspaces and increasing
biases. It is unchanged by independent positive dilations of the labels
and has diminishing increments under adjoining suppliers to any list.
\end{lemma}

\begin{proof}
For the empty list the root has value $s\dim E$ by homogeneity and
translation. The direct sum and merging rules pass through the derivative
(the root of a direct sum can itself be split at its label).
For $h>0$, raising the root bias from $0$ to $h$ amounts to adding
$(E,0,h)$. Its increment can only decrease when $L$ grows. Passing to
the right derivative proves monotonicity under adjoining a supplier.
An enlargement or a bias increase amounts to adjoining the improved
supplier.

Common dilation of the labels preserves rank by homogeneity. For
$0<t<1$, adjoin a shortened copy $u_t=(A,ty,t\theta)$ of an existing
supplier $u=(A,y,\theta)$. At every root bias $s\ge0$, the root
absorbs $u_0$ and the original supplier absorbs $u_1$. By convexity
along this segment and monotonicity under addition, $u_t$ also leaves
the value unchanged. Deleting $u$ can only decrease rank, so replacing
$u$ by $u_t$ can only decrease rank. Dilate the resulting list by
$1/t$ and shorten each of its other labels by $t$ to return to the
original list. This proves equality for one shortened label and hence
invariance under arbitrary independent positive dilations.

It remains to compare the rank gained by adjoining $u$ to two lists
$L\subseteq L'$. Shorten $u$ towards the root. The added rank stays
constant by dilation invariance, and in the limit it can be read from
the derivative in the bias of the shortened supplier.

To justify the derivative calculation, perturb the ordinary biases,
keeping shared occurrences identical. For almost every perturbation,
the root-augmented values for $L,L'$ and for their augmentations by
$u_{1/k}$, $k\ge1$, are all differentiable in their biases at root bias
zero. Indeed each finite convex value is differentiable almost
everywhere. Bias translation determines it from its restriction to
root bias zero, and scaling the additional bias by $1/k$ is an
invertible coordinate change. We may therefore make this countable
family of choices simultaneously.

At a differentiability point the sum of the ordinary suppliers' bias
derivatives equals the rank. Apply this to $L$ augmented by $u_{1/k}$,
differentiating in the actual bias of that supplier. As $k\to\infty$,
the derivatives at the old suppliers converge to those for $L$ alone.
In fact, as convex functions of the old biases the values converge
pointwise by continuity and redundancy of $u_0$. Left and right
difference quotients squeeze each partial derivative at a point of
differentiability of the limit. Independent dilation leaves the
augmented rank fixed. The bias derivative at $u_{1/k}$ therefore
converges to $r_E(L\cup\{u\})-r_E(L)$.
For every $k$ that derivative can only decrease on passing from $L$ to
$L'$, since raising the bias is equivalent to adding a supplier and
$\lambda$ has diminishing increments. This proves the desired comparison
at the perturbed biases. Finally approximate all ordinary biases
simultaneously from below. Monotonicity squeezes each perturbed rank
between its original value and its value after a common negative shift.
The latter converges to the original value by right continuity of the
right derivative in \eqref{flag:capdef}.
\end{proof}

When all suppliers lie at the tip, the threshold rank is
\[
 r_E(L-s)=\dim\!\sum_{\theta_i>s} A_i .
\]
After the shift by $s$, the root absorbs every supplier of nonpositive
bias as its own bias increases from zero. Independently dilating the
remaining labels to $(0,1)$ merges their subspaces into a single space
$A$. On $A$ the root is redundant for small positive root bias. On a
complement it is the sole supplier. This accounts for both the rank
formula and the strict inequality at a bias tie.

A single column illustrates what changes away from the tip. For the
exterior price, let $E=K$ and put the column at a site of radius $u>0$
with bias $\theta$. The two-supplier exterior pairing, computed below,
gives $\lambda_L(s)=s+u\log(1+e^{(\theta-s)/u})$. Thus
$r_E(L-s)=(1+e^{(s-\theta)/u})^{-1}$. At the bias tie the column
supplies half a dimension. To combine such fractional contributions
we need a dimension calculus that can also measure their overlaps.

\subsection{Dimension of subspace expressions}

Introduce a formal space $X$ of normalized dimension one and a symbol
$U_j$ for the subspace supplied by label $j$. The unrestricted label $(0,1)$
designates $X$ itself. At that label a supplier of $A$ has rank
$\dim A$ by the tip calculation. For a scalar matrix $B$ with labeled
columns write $r(B)$ for the threshold rank in its output space, using
one supplier per label with subspace the span of its columns.

Consider a variable $v_i\in U_{j_i}$ for each input column. Given
scalar matrices $P,Q$ over $K$ with those columns, form the expression
in $X^e$
\[
 T=\{Pv: Qv=0,\ v_i\in U_{j_i}\},
\]
where $e$ is the number of rows of $P$. It describes the outputs $Pv$
left after imposing the constraints $Qv=0$. Sums, intersections,
direct sums, and scalar images each admit such a description by
introducing variables and linear constraints. In module theory these
are positive primitive subspace expressions \cite{Prest}. An inclusion
between two expressions is \emph{universal} if it holds for every
substitution of an actual vector space for $X$ and actual subspaces for
the $U_j$, with the unrestricted label interpreted as the whole space.

For ordinary subspaces, projecting the joint image $(Pv,Qv)$ to $Qv$
has kernel isomorphic to $T$. The proposition shows that the difference
of the corresponding threshold ranks also defines a dimension and
satisfies the same rank--nullity identity.

\begin{proposition}[Dimension calculus]\label{flag:dimension}
The formula
\begin{equation}\label{flag:ppdim}
 \dim_*T=r\!\binom{P}{Q}-r(Q)
\end{equation}
is nonnegative, respects universal inclusions and adds on direct sums.
For any scalar map $M$,
\[
 \dim_*T=\dim_*M(T)+\dim_*(T\cap\ker M).
\]
In particular $r_E(L)=\dim_*\big(\sum_i A_i\otimes U_{j_i}\big)$, where
each summand denotes the image along a basis of $A_i$ of independent
variables with label $j_i$.
\end{proposition}

\begin{proof}
First consider adding unrestricted columns spanning a space $D$ to a
matrix $B$. The resulting rank is
\[
 \dim D+r(B\bmod D).
\]
Indeed, the unrestricted columns supply the full rank $\dim D$ on $D$.
Copies of those columns with any other label add zero rank there. By
invariance under unused ambient dimensions and diminishing increments
(Lemma~\ref{flag:capacity}), they still add zero rank in the presence of
$B$. Make these copies for every label of $B$. Within each label we can
now subtract components in $D$ from its columns while preserving their
span together with the copies. The remaining columns lie in a fixed
complement of $D$. Removing the copies and applying direct sum
additivity gives the formula.

For $B=\binom{P}{Q}$ take $D=K^e$ in the first output coordinates. The
gain from adding the unrestricted columns of $D$ is $e-\dim_*T$. It
lies in $[0,e]$ and can only decrease on adjoining more input columns.
Thus $\dim_*T$ is nonnegative and can only increase with more inputs.
Dropping constraint rows also increases this dimension. For each
dropped row adjoin an unrestricted column on its constraint axis to
the joint matrix. By the quotient formula, that axis contributes one
to both ranks in~\eqref{flag:ppdim}, leaving exactly the expression
with the row dropped.

Now suppose $T\subseteq T'$ universally, with respective data
$(P,Q)$ and $(P',Q')$. We need a scalar comparison between their
presentations. If $T$ has $s$ inputs, substitute
\[
 X=K^s/\operatorname{rowspan}(Q)
\]
and let $v_i$ be the classes of the standard basis vectors. For each
restricted label take $U_j$ to be spanned by its $v_i$. This choice
realizes $Qv=0$ and imposes exactly the row relations generated by
$Q$. Universal inclusion gives $v'$ of the allowed labels with
$P'v'=Pv$ and $Q'v'=0$. Expressing $v'$ in the generators $v_i$
produces a scalar matrix $A$. Each restricted row uses only columns
of its own label, and an unrestricted row may use any column. The
relations yield scalar matrices $J,J'$ with
\[
 P'A-P=JQ,\qquad Q'A=J'Q.
\]

Row operations therefore let us replace $P$ by $P'A$ and append the
redundant constraints $Q'Av=0$ in~\eqref{flag:ppdim}. Drop the old
constraints $Qv=0$, increasing dimension. The resulting joint matrix
is $\binom{P'}{Q'}A$. Each of its labeled columns lies in the span of
columns of $\binom{P'}{Q'}$ with that label and of unrestricted
columns. Copy the unrestricted columns into each label. These copies
contribute zero in both ranks by the quotient argument above. Enlarging
the column spans now gives $\dim_*T\le\dim_*T'$. Two presentations of
the same expression include each other universally, so their dimensions
agree.

Direct sums add by Lemma~\ref{flag:capacity}. The dimension of $M(T)$ is
$r\binom{MP}{Q}-r(Q)$. For $T\cap\ker M$ the added equations are $MPv=0$,
so the two dimension formulas telescope, since stacking $MP$ with $P,Q$
does not change their joint rank. Taking $Q$ to have no rows gives the
last assertion.
\end{proof}

Apply the kernel identity to addition on $U\oplus V$. Its kernel is a
scalar copy of $U\cap V$, so
\begin{equation}\label{flag:modular}
 \dim_* (U+V)+\dim_*(U\cap V)=\dim_*U+\dim_*V .
\end{equation}
We identify expressions $U,V$ if $U\cap V$ and $U+V$ have the same
dimension. For nested expressions of equal dimension, adjoining a
third expression by sum or intersection preserves equality of
dimension by \eqref{flag:modular}. The same holds for direct sums and
scalar images by Proposition~\ref{flag:dimension}. These observations,
applied through $U+V$, give transitivity and allow equivalent expressions
to be substituted in subsequent linear constructions.
Now specialize to the subspace symbols of a single site $y$. If
$\theta\le\theta'$, adding the lower-bias supplier to the higher one
leaves its rank unchanged. Thus $U_{y,\theta}+U_{y,\theta'}$ is
equivalent to $U_{y,\theta'}$. For any finite list of biases their
cumulative sums give universally nested representatives. Likewise,
in the pair of labels $(y,\theta),(ty,t\theta)$ we can dilate one
label until they agree and merge the pair. Its sum then has the same
dimension as each member. Hence $U_{ty,t\theta}$ and $U_{y,\theta}$
are equivalent for every $t>0$.

\subsection{The metric on flags}

An \emph{event} is a subspace expression in $X$, with its dimension
in $[0,1]$ interpreted as a probability. A finite flag $F$ is a
decreasing filtration of $X$ by such expressions, with finitely many
real jump levels. Its thresholds $F_{\ge a}$ range from $X$ to zero.
We take them to be universally nested, using cumulative sums when they
are initially nested only up to equivalence. Their dimensions are the
upper tails of a probability law on the levels. Two flags $F,G$ also
have a pair law, determined by
\begin{equation}\label{flag:pair}
 \Pr(F\ge a,G\ge b)=\dim_*(F_{\ge a}\cap G_{\ge b}).
\end{equation}
Indeed a cell mass is a successive difference of four intersection
dimensions. The intersection for the two weaker thresholds contains the
sum of the intersections obtained by strengthening either threshold.
Equation~\eqref{flag:modular} then gives nonnegativity.

A family of these pair laws may fail to admit a common probability
space. For example, take three distinct lines in a two-dimensional
space and normalize dimension by $1/2$. Each line has probability
$1/2$, while each pair has intersection probability zero. Three events
with these probabilities would have union probability $3/2$. We
therefore combine many events one pair at a time.

For a fixed pair, the grid of its jump levels has more structure.
For each upper set $S$ in this grid, let $U_S$ be the sum of the
intersections $F_{\ge a}\cap G_{\ge b}$ indexed by the minimal
elements $(a,b)$ of $S$. Unions correspond to sums and intersections
to intersections. These identities hold universally because for every
actual substitution the two filtrations admit a common adapted basis.
To construct one, work from the deeper grid intersections upwards.
Extend the bases already chosen for the two immediately deeper
intersections, whose common basis was fixed at the doubly deeper
intersection. Each $U_S$ is then a coordinate subspace in this basis.
Modularity and~\eqref{flag:pair} give
$\dim_*U_S=\Pr((F,G)\in S)$.

In particular, the upper level sets of a coordinatewise nondecreasing
function $f$ give a flag $f(F,G)$. Its law is the pushforward of the
pair law. The same grid gives the joint law of any two such functions
of $F,G$. This permits interpolation between two flags even when a
third flag has no joint law with them.

We use the notation
\[
 \langle F,G\rangle=\mathbb E[FG],\qquad
 \|F\|^2=\mathbb E[F^2],\qquad d(F,G)^2=\mathbb E[(F-G)^2],
\]
where each expectation uses the indicated law. The metric estimate
below will control what happens during insertion when one indicator is
changed. For the CAT(0) metric on filtrations of modular lattices of
finite length, see \cite[Proposition~2.6]{Cornut}.

\begin{proposition}[Flag metric]\label{flag:metric}
After identification at distance zero, finite flags have a CAT(0)
metric, with geodesic $H_t=((1-t)F+tG)$ defined by the two-flag
construction. The completion is Hadamard. Positive dilation is a
homothety. For fixed $Q$ the pairing with $Q$ is $\|Q\|$-Lipschitz.
\end{proposition}

\begin{proof}
The two-flag grid gives distances along the candidate geodesic
$H_t$. To compare it with a third flag $Q$ we first take
$Q=1_A$, the indicator flag of an event $A$. For each upper set $S$
of the $F,G$ grid put $w(S)=\dim_*(A\cap U_S)$. This function is
supermodular because the sum of two intersections with $A$ lies in
the intersection of $A$ with the sum. More explicitly, modularity
gives
\[
 w(S)+w(T)\le w(S\cup T)+w(S\cap T).
\]
Order the grid points $p_1,\ldots,p_s$ so each initial segment $S_i$
is upper, and set $S_0=\varnothing$. For coordinatewise nondecreasing
$f$, the expectation $\langle Q,f(F,G)\rangle$ equals
\[
 \min_{(p_i)}\ \sum_{i=1}^s f(p_i)\big(w(S_i)-w(S_{i-1})\big).
\]
Indeed the pair law gives equality for a decreasing order of $f$.
From any allowed order, swapping adjacent points inverted by $f$
preserves the upper-set condition, since those points are incomparable.
Supermodularity shows that such a swap can only decrease the sum.
The minimum is concave in $f$. Every finite $Q$ decomposes as a constant
plus a nonnegative combination of its threshold indicators, with the
corresponding equality of pairings. For every $Q$ this proves
\[
 \langle Q,H_t\rangle\ge(1-t)\langle Q,F\rangle+t\langle Q,G\rangle.
\]
Computing $\|H_t\|^2$ in the pair law and expanding now yields
\[
 d(Q,H_t)^2\le (1-t)d(Q,F)^2+t d(Q,G)^2-t(1-t)d(F,G)^2 .
\]
Nonnegativity of the left side, with
$t=d(Q,F)/(d(Q,F)+d(Q,G))$, gives the triangle inequality
(the zero-denominator case follows at $t=1/2$).
The joint law on the grid gives
$d(H_t,H_u)=|t-u|d(F,G)$, establishing the geodesic and CAT(0)
comparison. Dilation and homogeneity of the pairing follow from the
pair laws. Finally the triangle inequality gives
$|d(F,aQ)-d(G,aQ)|\le d(F,G)$ for $a>0$.
Taking the difference of squares and dividing by $2a$, then letting
$a\to\infty$, shows
$|\langle F,Q\rangle-\langle G,Q\rangle|\le\|Q\|d(F,G)$.
All these properties pass to the metric completion.
\end{proof}

The exterior application has strict thresholds at every real level.
Suppose a decreasing family $F_{>a}$ has dimensions equal to the
tails of a probability law of finite second moment. Truncate and
discretize the levels, using nested representatives at each finite
set of thresholds. Two such approximations can be coupled on their
combined chain. Their distance is the ordinary $L^2$ distance of the
two discretizations under the given law, so they converge in the flag
completion. For two families the pair upper tails are still the
dimensions of the threshold intersections. Indeed, by
\eqref{flag:modular} the change in an intersection dimension on
changing one threshold is bounded by the change in its marginal. The
upper tails are therefore right-continuous, and the continuous pairing
is the integral of the product under this pair law.

For each site $y$ take $F(y)_{>a}=U_{y,-a}$. Increasing $a$ lowers
the bias and gives a decreasing family. Whenever its tail dimensions
have finite second moment as above, it defines a flag $F(y)$. The tip
calculation gives $F(0)=0$, and equivalence of dilated labels gives
$F(ty)=tF(y)$ for $t>0$. For the exterior value, the single-column
calculation will identify the law at every site.

\subsection{Insertion along a column curve}

Write $F+G$ for twice the midpoint of $F,G$. On their two-flag grid,
this is the flag of the function $(f,g)\mapsto f+g$. Shifting its
arguments shifts this sum by the total shift. CAT(0) convexity makes
addition $1$-Lipschitz in each argument. Given events $H_1,\ldots,H_N$
in a fixed order, start with $V_0=0$ and form
$V_i=V_{i-1}+1_{H_i}$. At each step the grid of the current flag and the next
indicator defines the new flag, and their pair law gives its
distribution. The final flag
$V=V_N$ is the \emph{insertion sum} for this order.

Write $D_j$ for the current threshold at level $j\ge1$, and put
$D_0=X$. After inserting an event $H$, the threshold becomes
\begin{equation}\label{flag:thresholdupdate}
 D_j^{\mathrm{new}}=D_j+(H\cap D_{j-1})\qquad (j\ge1).
\end{equation}
Indeed, the upper set for $V_{i-1}+1_H\ge j$ consists of the old
threshold together with the part of level $j-1$ that meets $H$. Thus
after two insertions the first two thresholds are the sum and the
intersection of the events. In general $V$ is nonnegative and
integer-valued. Its first $m$ threshold dimensions sum to
$\mathbb E\min(m,V)$, the count that can fit into $m$ row directions.
The next lemma realizes this quantity as a lower bound on column rank.

\begin{lemma}[Insertion along a curve]\label{flag:insertion}
Let $v(t)\in K(t)^m$ be regular at zero and suppose its Taylor
coefficient columns span $K^m$. Read the coefficients in order of
degree, retaining each that increases the span. Denote the retained
basis by $e_1,\ldots,e_m$. Let $D_j$ be the thresholds obtained by
inserting $H_1,\ldots,H_N$ in that order. For every subspace expression
$U\subseteq X^m$ and $\varepsilon>0$, there are parameters with
\[
 \dim_*\!\left(U+\sum_{i=1}^N v(t_i)H_i\right)
 \ \ge\ \dim_*\!\left(U+\sum_{j=1}^m e_jD_j\right)-\varepsilon .
\]
One tuple $(t_1,\ldots,t_N)\in K^N$ can meet finitely many such tests
with the same curve and insertion order, allowing the events and $U$
to vary between tests. It can be chosen in any specified nonempty
principal open subset of $K^N$, avoiding poles and pairwise
coincidences. For $v(t)=(1,t,\ldots,t^{m-1})^\top$ and $U=0$, the
lower bound is $\mathbb E\min(m,V)-\varepsilon$.
\end{lemma}


For the curve $(1,t,t^2)^\top$, suppose the earlier events have
supplied $e_1L_1+e_2L_2+e_3L_3$ with $L_1\supseteq L_2\supseteq L_3$.
At $t=0$ the new column supplies $H$ along $e_1$. On $H\cap L_1$
the old space $e_1L_1$ can cancel the constant term of the new column.
Dividing by $t$ exposes $e_2$ at zero. On $H\cap L_2$ the old spaces
can cancel both lower terms, exposing $e_3$ after division by $t^2$.
Specializing all three expressions together captures the three
contributions in~\eqref{flag:thresholdupdate}.

\begin{proof}
We first record a specialization estimate for the dimension of $B(t)T$,
where $T$ is any fixed subspace expression and $B$ is a polynomial
matrix of degree at most $b$. At $q>b$ distinct scalars $a_i$,
\begin{equation}\label{flag:special}
 \sum_{i=1}^q\dim_* B(a_i)T\ \ge\ (q-b)\dim_* B(0)T .
\end{equation}
Take independent inputs from $T$ as the $q-b$ coefficients of a vector
polynomial $p$ of degree less than $q-b$. Evaluating $B(t)p(t)$ at the
$a_i$ places the joint output in $\bigoplus_i B(a_i)T$.
Scalar interpolation recovers its coefficients since its degree is less
than $q$. Its first $q-b$ coefficients are a lower triangular block map
of the inputs with diagonal block $B(0)$. Its image has dimension at
least $(q-b)\dim_*B(0)T$. Project onto the first output, whose
dimension is $\dim_*B(0)T$. The kernel contains the image obtained
by setting the first input to zero, a triangular map of the same
form with one fewer input. Proposition~\ref{flag:dimension} now
adds the bounds for the image and kernel. This proves
\eqref{flag:special} by induction. If $R=\dim_* B(0)T$, a set of
$q>b$ parameters each losing at least $\eta>0$ satisfies
$q(R-\eta)\ge(q-b)R$, whence $q\le bR/\eta$. In particular all
but finitely many $a$ have $\dim_* B(a)T\ge R-\eta$.
Clearing common denominators gives the same conclusion for rational
maps regular at zero.

Let $L_j$ be the thresholds before the last insertion and put $L_0=X$.
Consider the expression
\[
 U+v(t)H_N+\sum_{j=1}^m e_j L_j .
\]
Write $v(t)=\sum_h c_h t^h$ at zero and $e_j=c_{h_j}$. For each $j$
add an independent input from $H_N\cap L_{j-1}$ along
\[
 w_j(t)=t^{-h_j}\!\left(v(t)-\sum_{h<h_j}c_h t^h\right).
\]
This is regular at zero with $w_j(0)=e_j$. Every coefficient of the
subtracted prefix lies in the span of $e_1,\ldots,e_{j-1}$. Since
$L_{j-1}\subseteq L_i$ for $i<j$, the extra inputs leave the image
unchanged away from zero and the poles. At zero they supply all the
new terms $e_j(H_N\cap L_{j-1})$ in \eqref{flag:thresholdupdate}.
Apply the specialization estimate to the direct sum of all the fixed
input expressions, including the new inputs. At zero its image contains
$U+\sum_j e_jD_j$ for the updated thresholds. At every nonzero regular
parameter the image equals the expression before the extra inputs were
added. For each test all but finitely many $t_N$ therefore attain the
required bound with slack $\varepsilon/N$.

Choose $t_N$ good for every test and such that the prescribed principal
open has a nonempty slice in the remaining variables. For the latter
condition it suffices to keep a nonzero coefficient of its defining
polynomial, viewed as a polynomial in those remaining variables,
nonzero at $t_N$. There are only finitely many excluded values and $K$
is infinite. For $N=1$ this choice proves the result. For $N>1$ apply
induction to the earlier insertions with fixed part $U+v(t_N)H_N$ and
the remaining slack. Distinctness and pole
avoidance can be included in the initial principal open. The case $N=0$
is immediate.
\end{proof}
